% !PRINTED 636
\weberpage{636}
%<<CONT
dann ist nach \S.\,172, III. die Zerlegung von $p$ in Primfactoren im Körper $\Omega_m$ dieselbe, wie im Körper $\Omega_{m_1}$, und es ist
\[
\mathfrak{p}_n = \mathfrak{p}_{n+sm_1},
\]
wenn $s$ eine beliebige ganze Zahl ist. Wir setzen jetzt
\weq{18}{t = n + s m_1}
und lassen $n$ die Reihe der durch $q$ nicht theilbaren positiven ganzen Zahlen $<m_1$, und $s$ die Zahlenreihe $0,1,2,\ldots,m_2-1$ durchlaufen, dann durchläuft $t$ ein volles System durch $q$ nicht theilbarer Reste von $m$, und zwar das ganz bestimmte System, dessen Elemente alle $<m$ sind. Es ist dann nach \S.\,172
\weq{19}{p = \prod^{n} \mathfrak{p}_n.}

Ist ferner $t'$ eine Zahl, die der Congruenz
\weq{20}{t\,t' \equiv 1\ (\mathrm{mod}\ m)}
genügt, so ist, wenn $n'$ durch die Bedingung $n n' \equiv 1\ (\mathrm{mod}\ m_1)$ bestimmt ist,
\weq{21}{t' \equiv n'\ (\mathrm{mod}\ m_1).}

Danach lässt sich das Product
\[
\prod^{t} \mathfrak{p}_{t'}^{\,t}
\]
nach der Formel (16) bestimmen, die, auf den Körper $\Omega_{m_1}$ angewandt, der aus allen rationalen Functionen der $m_1^{\mathrm{ten}}$ Einheitswurzel $r_1 = r^{m_2}$ besteht, die Formel ergiebt:
\weq{22}{(r_1,\eta)^{m_1} = \prod^{n} \mathfrak{p}_{n'}^{\,n}.}

Es ist aber nach (18) und (21)
\[
\prod^{t} \mathfrak{p}_{t'}^{\,t} = \prod^{n} \mathfrak{p}_{n'}^{\,n m_2} \prod^{n'}\prod^{s} \mathfrak{p}_{n'}^{\,s m_1},
\]
und da $n'$ nach dem Modul $m_1$ dieselbe Zahlenreihe durchläuft wie $n$, so ist nach (19):
\[
\prod^{s}\prod^{n'} \mathfrak{p}_{n'}^{\,s m_1} = p^{m_1 \Sigma s} = p^{\frac{1}{2}m(m_2-1)},
\]
und es ergiebt sich:
\[
\prod^{t} \mathfrak{p}_{t'}^{\,t} = \left[p^{\frac{1}{2}(m_2-1)}\,(r^{m_2},\eta)\right]^m.
\]

Setzen wir
\weq{23}{\varrho = p^{\frac{1}{2}(m_2-1)}\,(r_1,\eta), \qquad \varrho_n = p^{\frac{1}{2}(m_2-1)}\,(r_1^{\,n},\eta),}
so ist $\varrho$ eine \gesperrt{Kreistheilungszahl}, weil ja auch $\sqrt{p}$ als Werth
%CONT>>
